\specsection{测度的逆极限}

这一理论主要是为满足概率论的需要而发展起来的。关于随机变量有限序列 \(X_{1},\ldots,X_{n}\) 的问题，原则上在知道该序列的律 \(P_{X}\) 时就解决了：这是 \(R^{n}\) 上质量为 1 的一个正测度，使得同时得到不等式 \(a_{1}\leqslant X_{1}\leqslant b_{1},\ldots,a_{n}\leqslant X_{n}\leqslant b_{n}\) 的概率等于 \(\mathrm{P}_{\mathrm{X}}(\mathrm{C})\)，其中 C 是闭方体 \([a_{1},b_{1}]\times\cdots\times[a_{n},b_{n}]\)（在 \(R^{n}\) 中）。在实践中，测度 \(P_{X}\) 或者有离散支撑，或者关于 Lebesgue 测度有密度。当处理随机变量的无穷序列 \((\mathrm{X}_{n})_{n\geqslant1}\) 时，律 \(P_{n}\)（部分序列 \((\mathrm{X}_{1},\ldots,\mathrm{X}_{n})\) 的）对每个整数 \(n\geqslant1\) 都是已知的；这些数据满足一个相容性条件，它表示序列 \((\mathrm{P}_{n})_{n\geqslant1}\) 是测度的一个逆系统（或“射影系统”）。直到大约 1920 年，涉及无穷序列的事件的概率还是以或多或少隐含的方式，由有限情形概率的“自然”极限过渡来定义的；例如，人们假设一局赌博终止的概率是“它至多在 n 步内终止”的概率当 n 趋于无穷时的极限。自然，这样的理论不太协调，而且没有什么能排除“悖论”的出现：同一个概率，按照人们用两个（各自都同样“自然的”）程序中的这一个或那一个来计算，会得到两个不同的估计。

Steinhaus (V) 似乎是最早感到有必要（对掷硬币游戏）不仅考虑逆系统 \((\mathbf{P}_{n})_{n\geqslant1}\)、而且考虑其极限的人。稍早一些，在 1919 年，Daniell (VI, b)) 已经一般地证明了这种逆极限的存在，\(^{(7)}\) 但这一结果在欧洲似乎一直不为人所知。它在 1933 年被 Kolmogoroff 在著作 (XII) 中重新发现，作者在该著作中阐述了概率论的公理化概念。Daniell 与 Kolmogoroff 的证明利用了一个紧性论证，它实质上与我们在 §4, No.~3 的 Th. 2 中所用的相同，并且基于 Dini 定理。

Daniell--Kolmogoroff 定理对随机序列 \((\mathbf{X}_{n})_{n\geqslant1}\) 的情形已无可挑剔，但从 1935 年起由 Kolmogoroff、Feller 与 Doob 开展的随机函数的研究则隐含着完全另一种层次的困难。例如考虑 R 的一个区间 T，表示一个“随机过程”的观察时刻所成的集合；可能轨迹所成的集合是乘积空间 \(R^{T}\)，它被视为部分乘积 \(R^{H}\) 的逆极限，其中 H 取遍 T 的有限子集所成的集合；一般假设给定了测度的一个逆系统 \((\mu_{\mathrm{H}})\)（cf.~§4, No.~2）。Kolmogoroff 定理确实给出了 \(R^{T}\) 上的一个测度，但它定义在一个比 Borel σ-集环小得多的 σ-集环上。\(^{(8)}\) Kolmogoroff 构造的一个变体（它给出拓扑空间上的一个测度）应归功于 Kakutani（Proc. Imp. Acad. Tokyo 19 (1943), 184--188），此后被多次重新发现：把 \(\mu_{H}\) 视为 \(\overline{R}^{H}\) 上由 \(R^{H}\) 支撑的测度；\(^{(9)}\) 紧空间 \(E = \overline{R}^{T}\) 是有限乘积 \(\overline{R}^{H}\) 的逆极限，于是可以把 E 上的测度 \(\mu\) 定义为诸 \(\mu_{H}\) 的逆极限（cf.~Ch. III, §4, No.~5）。然而，这一做法有一个严重的不便；\(\overline{R}^{T}\) 的元素不具备任何能使过程的概率研究取得进展的正则性质---甚至不能简单地消去寄生值 \(\pm\infty\)（由 R 的紧化 \(\overline{R}\) 所引入的）。补救的办法是把测度 \(\mu\)（\(\overline{R}^{T}\) 上的）诱导到一个特定的子空间上（例如 Brown 运动情形中的 \(\mathcal{C}(T)\)）；基本的困难来自这样一个事实：一个函数空间，即使是通常类型的，也不一定是 \(\mu\) -可测的（在 \(\overline{R}^{T}\) 中），甚至连函数空间的选择都可能成问题。\(^{(10)}\)

决定性的一步是 Prokhorov 于 1956 年在一部对随机过程理论产生了决定性影响的著作 (XIII) 中迈出的。通过把 Wiener 在上面分析的文章中所用的方法放进一个适当的公理化形式，他建立了函数空间上测度逆极限的一个一般存在性定理，它是对应于波兰空间的 §4, No.~2 的 Th. 1 的一个特殊情形。

更受限的一类逆系统由 Bochner (XIV) 于 1947 年引入；这是由有限维实向量空间与满射线性映射组成的逆系统。这样一个系统的逆极限可以以自然的方式等同于代数对偶 \(\mathbf{E}^*\)（实向量空间 \(\mathbf{E}\) 的），赋以弱拓扑 \(\sigma (\mathbf{E}^{*},\mathbf{E})\)；相应的测度逆系统有一个极限，即测度 \(\mu\)（定义在一个比 \(\mathbf{E}^*\) 的 Borel σ-集环小得多的 σ-集环上）。Bochner 通过 Fourier 变换完全刻画了这种“预测度”，Fourier 变换是 \(\mathbf{E}\) 上的一个函数。但在 \(\mathbf{E}\) 上没有拓扑时，这一结果几乎无法使用，这时就必须考察把 \(\mu\) 视为拓扑对偶 \(\mathbf{E}'\)（\(\mathbf{E}\) 的）上的测度的可能性。以独立的方式，R. Fortet 与 E. Mourier 在寻求把概率论的某些经典结果（大数定律、中心极限定理）推广到取值于 Banach 空间的随机变量时，也揭示了 Fourier 变换在这类问题中所起的基本作用。但实质性的进展直到 1956 年才到来：Gelfand (XV, \(b\) ) 提出，Fourier 变换的自然框架不是 Banach 空间或 Hilbert 空间，而是核 Fréchet 空间。他猜想这样的空间上每个正型的连续函数是其对偶上某个测度的 Fourier 变换，这一结果不久由 Minlos (XVI) 证明。它的重要性首先在于它适用于广义函数空间，而且几乎全部的函数空间都是广义函数空间的 Borel 子集（因而它构成一个比 \(\mathbf{R}^{\mathrm{T}}\) 好得多的容器）。\(^{(11)}\) 随机广义函数理论是一个蓬勃发展的领域，我们只需请读者参考 Gelfand 与 Vilenkin 的书 (XVII)。

刚才提到的关于逆极限的结果利用了基空间上拓扑的存在性。人们可能会问，在“抽象”测度的情形是否存在类似的理论。Von Neumann 于 1935 年证明了乘积测度在所有情形下的存在性，但 Jessen 与 Andersen (XVIII) 发现的一个反例使每个测度逆系统都容许极限的希望破灭了。人们发现了两种权宜之计：1949 年，C. Ionescu-Tulcea 借助适当分解的存在性建立了可数逆极限的存在，\(^{(12)}\) 这是研究 Markoff 过程的一个非常有趣的结果；此外，人们已经观察到，空间的拓扑只是通过紧子集所成的集合这一中介介入。因此，很自然地试图在抽象理论内部、借助集合的紧类概念把这种情形公理化。这项工作于 1953 年由 Marczewski（他用这种方法建立了一个抽象逆极限定理）与 Ryll-Nardzewski（他处理了测度的分解）完成。\(^{(13)}\)

